\section{Density, Radius, and Regularization Methods}
\subsection{Radial Density}
Radial density defined as $\rho(r) = \langle \Psi(r')| \delta(r - r') | \Psi(r')\rangle$. Let $\hat{\rho}(r) = \delta(r-r')$. In our case the density comes out to be $\rho(r)=\Gamma u(r)^2$. In more complicated systems, the result is more complicated. 

\subsubsection{Complex Scaling the Density Operator}
The $\delta(r - r')$ is not something that can be complex-rotated in a traditional sense. It's a distribution, and it definitely doesn't behave in the well-defined way that most observables do (the conditions outlined in \cref{sssection:obs_phi_indep}. As such, there's not really a well-defined way to define the rotated version of $\delta(r - r')$. Instead, we can define an operator $\rho_\theta(r) = (\psi_\theta|\delta(r-r')|\psi_\theta)$. I.e., we don't attempt to rotate the Dirac delta. This expression is of course $\theta$-dependent, and so the goal would be to compute it at $\theta=0$. As such, we need to involve back-rotation. 

In our case, $\rho_\theta(r) = (\psi_\theta|\delta(r-r')|\psi_\theta) =\Gamma e^{i\theta}u_\theta(r)^2$. This varies with $\theta$, but the interior stays nearly constant. The $e^{i\theta}$ factor comes from the $e^{3i\theta/2}$ factor in the 3D complex scaling operation $\Psi_\theta(r) = U_\theta\Psi(r) = e^{i3\theta/2}\Psi(re^{i\theta})$.
\subsubsection{Finding the density through backrotation}
The back-rotated density is just $\rho^{\text{back}}(r) = \Gamma u_0(r)^2 =  \Gamma u_\theta(re^{-i\theta})^2$ because the back-rotated wavefunction is just the original Gamow state. Casting this with respect to the psuedo-rotated density:

$$
\rho_{\theta, BR}(r) = e^{-i\theta} \rho_\theta(re^{-i\theta}),
$$

One can alternatively show this with a similar integral argument that justified $\rho_\theta$ by back-rotating the resonance wavefunction: $\Psi_{\theta,BR}(r) = e^{-i3\theta/2}\frac{u_\theta(e^{-i\theta}r)}{e^{-i\theta}r}$ to recover the Gamow state, then computing the density with the standard delta function.
\subsection{Radius squared operator ($r^2$)}
\subsubsection{Rotated radius squared operator}
$\langle r^2\rangle_\phi$ is radius squared operator of CS'd state $\psi_\phi$. $r\rightarrow e^{i\phi} r$, so
$$
\Psi_\theta(r) = U_\theta \Psi(r) = e^{3i\theta/2}\Psi(re^{i\theta})
$$
$$
= e^{i\theta/2}\frac{u(re^{i\theta})}{r}Y^0_0.
$$
The $e^{3i\theta/2}$ factor comes from requiring that the complex scaling operator is unitary in 3D.
$$
\langle r^2_\theta \rangle =\frac{1}{4}\frac{(\Psi_\theta|r^2_\theta|\Psi_\theta)}{(\Psi_\theta|\Psi_\theta)}
$$
$$
=\frac{1}{4}\frac{e^{i\theta}\int \;d\Omega\; (Y_0^0(\Omega))^2 \int\; dr \; r^2e^{2i\theta} u_\theta(r)^2}{e^{i\theta}\int \;d\Omega\; (Y^0_0(\Omega))^2\int\;dr \; u_\theta(r)^2} = \frac{e^{2i\theta}}{4} \int\;dr\; r^2 u_\theta(r)^2
$$

the $1/4$ factor is just to average out the CM. $\Gamma$ factor is angular contribution (we work exclusively in S wave) and is ignored (though here it cancels out anyway). DVR code is working in relative coordinate of reduced radial equation, so $u_\phi$ are the eigenvectors.

In DVR quadrature:
$$
\langle r^2_\theta\rangle = e^{2i\theta}\frac{1}{4}\sum_k r_k^2 u_\theta(r_k)^2 \omega_k = e^{2i\theta} \frac{1}{4} \sum_k c_k^2 r_k^2,
$$
where $c_k=u_\theta(r_k)\sqrt{\omega_k}$ is the $k$-th DVR expansion coefficient, and the quadrature weights $\omega_k$ cancel due to c-normalization.

\subsubsection{Back-rotated radius squared operator}
$\langle r^2 \rangle_{\text{BR}}$ is the expectation value of the back-rotated Gamow state. Here $r^2$ is un-rotated, and $\Psi_{\theta=0}=U_\theta^{-1}\Psi_{\theta}=e^{-3i\theta/2}\Psi_{\theta}(re^{-i\phi}) =\frac{u_\theta(e^{-i\theta} r)}{r}Y^0_0$. Thus,

$$
\langle r^2 \rangle_{\text{BR}} = \frac{1}{4}\frac{\int \;d\Omega\;(Y_0^0(\Omega))^2\int dr\;r^2 u_\phi^2 (e^{-i\phi}r)}{\int\;d\Omega\;(Y_0^0(\Omega))^2\int dr\;u_\phi^2 (e^{-i\phi}r)} = \frac{1}{4} \frac{\int dr\;r^2 u_\phi^2 (e^{-i\phi}r)}{\int dr\;u_\phi^2 (e^{-i\phi}r)}
$$
$$
=\frac{1}{4}\frac{e^{i3\theta}\lim_{\epsilon\rightarrow \infty}\int_0^{\epsilon e^{-i\theta}}\;dr'\;r'^2 u_\phi(r')}{e^{i\theta}\lim_{\epsilon\rightarrow \infty}\int_0^{\epsilon e^{-i\theta}}\;dr'\;u_\phi(r')}= \frac{e^{2i\theta}}{4}\frac{\lim_{\epsilon\rightarrow \infty}\int_0^{\epsilon e^{-i\theta}}\;dr'\;r'^2 u_\phi(r')}{\lim_{\epsilon\rightarrow \infty}\int_0^{\epsilon e^{-i\theta}}\;dr'\;u_\phi(r')},
$$

where $u_\phi(re^{-i\theta}) = u(r)$ is the original Gamow state, and $r'=re^{-i\theta}$. The integral bounds for this new variable are $r'$ $\rightarrow 0$ to $r'\rightarrow\infty$, but, formally, the integral is to be taken along the ray along $e^{-i\theta}$, so $r'\rightarrow 0$ to $r'\rightarrow \infty e^{-i\theta}$. Thus, the ratio above is formally $\infty/\infty$, but it is well-defined as the boundary value $\theta\rightarrow0$ of 
$$
f(\theta)=\frac{e^{2i\theta}\int^\infty_0\;dr\;u_\theta(r)^2r^2}{\int_0^\infty\;dr\; u_\theta(r)^2}
$$
which is constant in the "good" $\theta$ sector by the ABC theorem. I.e., this singularity in the function $\langle r^2_\theta\rangle(\theta)$, while not formally defined, approaches a finite limit. Formally, this ratio is only defined in the analytic continuation of this expectation value. Since the ratio is the same constant through the wedge, the analytic continuation to $\theta=0$ will be the same constant by the identity theorem.

In DVR quadrature,
$$
\langle r^2 \rangle_{\text{BR}} = \frac{1}{4}\frac{\sum_k r_k^2 u_\phi^2 (e^{-i\phi}r_k) \omega_k}{\sum_k u_\theta(e^{-i\theta}r_k)^2\omega_k}.
$$
We normalize the states wrt the DVR basis, so the denominator is 1. The back-rotated radius computed in this way is ill-posed as errors accumulate due to back-rotating the decaying state back to the divergent Gamow state.
\subsection{Tikhonov Regularization}

$$
\langle r^2 \rangle_{\text{BR}} = \int dr\; r^2 u_\theta(e^{-i\theta}r)e^{2i\theta}
$$
Let $f_\theta(r) = r^2 u^2_\theta(r)$. Then, 
$$
\langle r^2 \rangle_{\text{BR}} = e^{4i\theta}\int dr\; f_\theta(e^{-i\theta} r).
$$
Changing variables $x=-\ln(r/r_0)$ gives $\tilde{f}_\theta(x)=f_\theta(r_0e^{-x})$ (domain goes from $(0, \infty)\rightarrow (-\infty,\infty)$. Fourier transform:
$$
\hat{f}_\theta(\xi) = \frac{1}{\sqrt{2\pi}} \int_{-\infty}^\infty dx\; e^{-ix\xi}\tilde{f}_\theta(x)
$$
Now, analytically continuing to complex plane: $\tilde{f}_\theta(x+iy)=f_\theta(r_0e^{-x}e^{-iy})$. So, $\tilde{f}_\theta(x+i\theta) = f_\theta(e^{-i\theta}r)$.
$$
\tilde{f}_\theta(x + i\theta) = \frac{1}{\sqrt{2\pi}} \int_{-\infty}^\infty d\xi\; e^{i(x+i\theta)\xi}\hat{f}_\theta(\xi)
$$
$$
= \frac{1}{\sqrt{2\pi}} \int_{-\infty}^\infty d\xi\; e^{ix\xi}e^{-\theta\xi}\hat{f}_\theta(\xi).
$$
Tikhonov cuts off the UV modes: $\Gamma_{\text{Tikhonov}} = \frac{1}{1+\kappa e^{-2\theta\xi}}$. Thus,
$$
\tilde{f}_{\text{reg},\theta}(x + i\theta)=  \frac{1}{\sqrt{2\pi}} \int_{-\infty}^\infty d\xi\; e^{ix\xi}e^{-\theta\xi}\hat{f}_\theta(\xi)\Gamma_{\text{Tikhonov}}.
$$

Thus,
$$
\langle r^2\rangle_{\text{BR}} = e^{4i\theta} \int_0^\infty dr \;\tilde{f}_{\text{reg},\theta}(r) = e^{4i\theta}\int_\infty^{-\infty}-r_0e^{-x}\;dx\;\tilde{f}_{\text{reg},\theta}(x+i\theta)
$$
$$
= e^{4i\theta}\int_{-\infty}^{\infty}r_0e^{-x}\;dx\;\tilde{f}_{\text{reg},\theta}(x+i\theta)
$$
\subsection{Zeldovich Regularization}
To be added.
